% Curated notation list: page references are generated by LaTeX.
\chapter*{Index of Notation}
\addcontentsline{toc}{chapter}{Index of Notation}
\markboth{Index of Notation}{Index of Notation}
This grouped list records the working convention and gives the page of a
definition or principal discussion. It is a notation guide, not a list of
every occurrence. Local hypotheses remain those of the referenced text.

\paragraph{Historical conventions.}
Some inherited passages use $N_0$ for the present $N$, $N$ for $J-N$,
$M_0$ for the present $M$, and $M$ for $M-\tau\omega$. These are historical
aliases, not additional operators. The ongoing notation revision removes
those aliases chapter by chapter. In the partial-circulant collision-spectrum
section, the locally used $N_0=J-N$ is a different alias and should be read
only with its local definition.

\begingroup
\newcommand{\notationentry}[3]{\item[$#1$] #2\quad\textit{p.~\pageref{#3}.}}
\section*{Graphs and transformations}
\begin{description}
\notationentry{\mathcal A}{Adjacency matrix; distinguish from the stochastic normalization and the locally reused family symbol in the partial-circulant chapter.}{defn:1}
\notationentry{R,B}{Letter transformations, represented by binary stochastic matrices.}{defn:1}
\notationentry{A}{Normalized transition matrix; in the two-letter equal-weight setting, $(R+B)/2$.}{sec:6.1}
\notationentry{\Delta}{Letter-difference matrix $(R-B)/2$.}{sec:6.1}
\notationentry{u}{All-ones row, of the dimension required by the representation level.}{sec:6.1}
\notationentry{J}{All-ones matrix $u^{\mathsf T}u$.}{sec:6.1}
\end{description}
\section*{Semigroups, kernel partitions, and ranges}
\begin{description}
\notationentry{\mathcal S}{Transformation semigroup generated by the letters.}{sec:6.1}
\notationentry{\mathcal K}{Minimum-rank ideal (semigroup kernel); not the nullspace of a matrix.}{sec:6.1}
\notationentry{r}{Kernel rank, the minimum transformation rank; not generally the rank of the averaged collision matrix.}{eproj:section}
\notationentry{e_{ij}}{Distinguished idempotent with kernel partition $i$ and range $j$.}{eproj:section}
\notationentry{\mathcal P_i,\mathcal R_j}{Kernel partition and kernel range; the latter is a transversal of each kernel partition.}{eproj:section}
\notationentry{\rho(k)}{Range-indicator row of a kernel transformation; shape $1	imes n$.}{sec:6.1}
\notationentry{\rho_j}{Diagonal indicator matrix of range $\mathcal R_j$; shape $n	imes n$.}{notation:range-projections}
\end{description}
\section*{Probability and stationary data}
\begin{description}
\notationentry{\lambda}{Limiting kernel probability measure used in kernel expectations; use the appropriate Cesaro convention.}{eproj:section}
\notationentry{\alpha_i,\beta_j}{Partition and range marginals of the kernel measure.}{eproj:section}
\notationentry{\pi}{Normalized stationary row of $A$.}{sec:6.1}
\notationentry{\Omega}{Stationary projection $u^{\mathsf T}\pi$ in the irreducible setting.}{sec:6.1}
\notationentry{\omega}{Diagonal matrix $\operatorname{diag}(\pi)$.}{sec:6.1}
\end{description}
\section*{Kernel expectations and projections}
\begin{description}
\notationentry{N}{Collision expectation $\mathbb E_\lambda[KK^{\mathsf T}]$; symmetric $n	imes n$ matrix.}{notation:expectations}
\notationentry{M}{Weighted range expectation $\mathbb E_\lambda[K^{\mathsf T}JK]$; symmetric $n	imes n$ matrix.}{notation:expectations}
\notationentry{J-N}{Pair-survival matrix, complementary to the collision expectation.}{notation:expectations}
\notationentry{\tau}{Trace $\operatorname{tr}M$ in the kernel-operator convention; the potential-operator section also uses this letter as an interpolation parameter.}{sec:6.3}
\notationentry{M-\tau\omega}{Centered range expectation; distinguish from $M$.}{sec:6.3}
\notationentry{\widetilde M}{Pair-range observable $\mathbb E[K^{\mathsf T}(J-I)K]$; distinct from the centered operator $M-\tau\omega$.}{propMN}
\notationentry{C_i}{Partition coincidence matrix $e_{ij}e_{ij}^{\mathsf T}$; independent of the range index $j$.}{notation:expectations}
\notationentry{P_j}{Range-fixed projection $N
ho_j$, acting on row vectors.}{notation:range-projections}
\notationentry{Q_i}{Partition-fixed projection; $rC_i\omega$ under the stated equipartition hypotheses.}{notation:partition-projections}
\notationentry{\overline e}{Product-marginal average of distinguished idempotents; not generally an idempotent.}{notation:average-idempotents}
\end{description}
\section*{Higher representations and inclusion maps}
\begin{description}
\notationentry{f^{\otimes k}}{Common-transformation action on the $k$th tensor power.}{hierarchy:collision-quotient}
\notationentry{\operatorname{Sym}^k(V)}{Symmetric power: unordered configurations with multiplicity.}{hierarchy:collision-quotient}
\notationentry{F_\ell}{Collision filtration by at most $\ell$ occupied vertices, at fixed degree.}{hierarchy:collision-quotient}
\notationentry{f^{\vee k}}{Degree-$k$ zeon action: a $k$-subset maps to its image or to zero after a collision.}{hierarchy:collision-quotient}
\notationentry{A_k}{Common-letter average at degree $k$; distinguish it from an independent-particle power of $A$.}{hierarchy:collision-quotient}
\notationentry{\Omega_k}{Higher-level limiting projection; interpreted through the stated Abel or Cesaro limit.}{sec:7.3}
\notationentry{u_\ell}{Probability column that an $\ell$-subset is split by a kernel partition.}{organized:inclusion-fields}
\notationentry{\eta_\ell}{Normalized range-subset probability row; distinguish from inherited higher-level row normalizations.}{organized:inclusion-fields}
\notationentry{E^{\ell,k}}{Inclusion matrix from $\ell$-subsets to $k$-subsets.}{organized:inclusion-fields}
\notationentry{D_{\ell,k}}{Normalized inclusion matrix $\binom{\ell}{k}^{-1}E^{\ell,k}$.}{organized:inclusion-fields}
\end{description}
\section*{Rank tests and potential operators}
\begin{description}
\notationentry{\Sigma(v,T)}{Cycle matrix with rows $vT,\ldots,vT^d$, where $d=\operatorname{rank}T$.}{defn:6.1}
\notationentry{Z}{Fundamental inverse $(I-A+\Omega)^{-1}$.}{potential:section}
\notationentry{V}{Potential operator $\Delta Z$; elsewhere $V$ may denote a vector space or vertex set.}{potential:section}
\notationentry{A_\tau}{Biased transition matrix $A+	au\Delta$; here $	au$ is an interpolation parameter, not $\operatorname{tr}M$.}{potential:interpolation}
\end{description}
\section*{Partial-circulant local notation}
\begin{description}
\notationentry{M_i}{Partial-circulant generator: $1M_i=i$ and $xM_i=x-1$ for $x>1$; not the expectation operator $M$.}{notation:ptc-definitions}
\notationentry{m,d,C}{Maximum generator index, gcd of generator indices, and recurrent core $\{1,\ldots,m\}$.}{notation:ptc-definitions}
\notationentry{P_a,C_a}{Residue block modulo $d$ and its intersection with the core (printed as $P_r,C_r$ with a local residue index $r$); not the range projection $P_j$.}{notation:ptc-definitions}
\end{description}
\endgroup
